\chapter{摘\texorpdfstring{\quad}{}要}

随着全球产业变革的推进，机器人的应用正从结构化的工业环境转向高度非结构化的居家场景。然而，现有的机器人模仿学习系统仍面对专家示教数据匮乏、非结构化环境约束复杂以及现实多模态任务决策困难等挑战。因此，研究基于观测输入的机器人模仿学习方法，对于推动机器人智能化水平的提升具有重要意义。

首先，针对模仿学习系统中专家数据稀缺的难题，本文基于机器人书法任务提出一种基于图像观测输入的书法技能模仿学习框架，实现从海量汉字书法图像中学习书法技能。该框架利用八邻域优先级算法实现笔画拓扑骨架的精准分割，并结合模糊高斯混合模型解决高曲率笔画的非线性拟合难题。通过引入偏差分析算法，从像素中重构表征笔画粗细变化的协方差序列，并结合动态运动原语实现了书法技能在不同空间尺度下的自适应泛化。搭建了基于艾利特机器人的汉字书写实验平台，在六个书写评价指标的平均提升了10.16\%，并测试了多个不同汉字的书写。

其次，针对模仿学习系统在非结构化环境中难以适应异构约束的问题，本文提出了基于异构约束观测输入的模仿学习网络。该网络通过编码空间障碍物位置或运动学限制等约束观测信息，利用零卷积初始化与残差注入机制，在冻结的扩散策略骨干中引入引导偏置，将物理约束信息注入到策略网络。仿真实验论证了不同物理约束在特征空间中的解耦特性，实现了多约束条件的零样本组合与即插即用式适配。在具有障碍物的仿真任务中，该方法达到了 90.5\% 的有效成功率，大幅增强了机器人复杂操作的安全性。

最后，为了提升模仿学习系统在现实环境中处理复合约束，动态约束和多任务泛化的能力，本文构建了基于多模态观测输入的模仿学习系统。利用自监督学习策略构建语音去噪网络，实现了基于语音指令的约束模块动态调度，在真实避障场景、运动学约束场景和复合约束场景中的成功率均在80\%以上。集成了实时目标检测模块提取障碍物几何特征，构建了闭环感知反馈回路，在长程动态避障分拣场景中的有效成功率达到 100\%。设计了基于对比学习的潜动作编码器，构建了多构型统一的动作表征，在后训练实现一倍的收敛速度提升。通过引入动作后处理机制，解决了推理频率与执行频率不匹配导致的运动抖动，使叠毛巾任务的成功率提升至 72\%，抓取任务成功率达到 90\%，提升了机器人在执行复杂任务时的运动稳定性与多任务泛化能力。

\keywordsCN{机器人模仿学习；高斯混合模型；动态运动原语；扩散策略；视觉-语言-动作模型}

\chapter{Abstract}
With the advancement of global industrial transformation, robots are progressively transitioning from structured industrial environments to highly unstructured domestic service scenarios. Nevertheless, existing robot imitation learning systems continue to face significant challenges, including the scarcity of expert demonstration data, the complexity of unstructured environmental constraints, and the difficulty of multi-modal task decision-making in real-world scenarios. Therefore, research on robot imitation learning methods based on observational inputs is of considerable significance for advancing robot intelligence and facilitating the large-scale deployment of embodied intelligence.

Firstly, to address the challenge of scarce expert data in imitation learning systems, this dissertation proposes an image observation-based calligraphy skill imitation learning framework leveraging robot calligraphy tasks, enabling skill acquisition from large-scale Chinese character calligraphy images. The framework employs an eight-neighbor priority algorithm to achieve precise segmentation of stroke topological skeletons, and incorporates a fuzzy Gaussian mixture model to address the nonlinear fitting challenges inherent in high-curvature strokes. Furthermore, a deviation analysis algorithm is introduced to reconstruct the covariance sequence characterizing stroke width variations from pixel data. This covariance information is then integrated with dynamic movement primitives to achieve adaptive generalization of calligraphy skills across different spatial scales. An experimental platform based on an Elite robot is established for Chinese character writing tasks. The proposed framework demonstrates an average improvement of 10.16\% across six writing evaluation metrics, and successful generalization is validated through testing on multiple different Chinese characters.

Secondly, to address the problem of imitation learning systems struggling to adapt to heterogeneous constraints in unstructured environments, this dissertation proposes an imitation learning network based on heterogeneous constraint observation inputs. The network encodes constraint observation information such as spatial obstacle positions or kinematic limitations, and utilizes zero- convolution initialization with residual injection mechanisms to introduce guiding biases into a frozen diffusion policy backbone, thereby injecting physical constraint information into the policy network. Simulation experiments demonstrate the decoupling characteristics of different physical constraints in the feature space, achieving zero-shot composition and plug-and-play adaptation of multiple constraint conditions. In obstacle-avoidance simulation tasks, this method achieves an effective success rate of 90.5\%, significantly enhancing robot operation safety in complex scenarios.

Finally, to enhance the capability of imitation learning systems to handle composite constraints, dynamic constraints, and multi-task generalization in real-world environments, this dissertation constructs an imitation learning system based on multi-modal observation inputs. A speech denoising network is built using a self-supervised learning strategy to achieve dynamic scheduling of constraint modules based on voice commands, with success rates exceeding 80\% in real obstacle-avoidance scenarios, kinematic constraint scenarios, and composite constraint scenarios. A real-time object detection module is integrated to extract obstacle geometric features and construct a closed-loop perceptual feedback loop, achieving an effective success rate of 100\% in long-horizon dynamic obstacle-avoidance sorting tasks. A latent action encoder based on contrastive learning is designed to construct a unified action representation across multiple configurations, achieving a two-fold improvement in convergence speed during post-training. By introducing an action post-processing mechanism, the motion jitter caused by the mismatch between inference frequency and execution frequency is resolved, raising the success rate of towel-folding tasks to 72\% and achieving a 90\% success rate in grasping tasks, thereby enhancing motion stability and multi-task generalization capabilities of robots when executing complex tasks.

\keywordsEN{Robotic Imitation Learning; Gaussian Mixture Model; Dynamic Movement Primitives; Diffusion Policy; Vision-Language-Action Models}